\subsection{\problemblock{*Cg} Data Block}\label{sec:block-cg}

One assumption in TAOS is that a control system is capable of maintaining the
vehicle body attitude. For example, a constant-angle-of-attack guidance scheme
assumes that a control system can maintain that attitude. On many aerospace
vehicles the control system uses aerodynamic devices such as flaps, canards, and
elevators. The settings required for steady flight at a given body attitude depend on
the vehicle center of gravity; maintaining a constant angle of attack can require
different control settings as the center of gravity changes.

The aerodynamic coefficients used by TAOS represent steady, trimmed flight. They
are therefore often functions of center of gravity. Center of gravity, in turn, is often a
function of vehicle weight. As fuel or propellant is consumed, weight and center of
gravity change, and this affects the aerodynamics.

TAOS handles this relationship by allowing aerodynamic tables to be functions of
the \statevar{cg} state variable and by defining \statevar{cg} with a
\problemblock{*cg} block. The value may be a constant or may be computed from a
table. The block is optional; center of gravity defaults to zero when no block is
present.

The following block sets center of gravity to a constant value:

\begin{lstlisting}[style=taosmanual]
*cg cg=0.69
\end{lstlisting}

Center of gravity can be expressed in inches, as percent of body length, or in any
other convenient convention, provided the same convention is used in both the
\problemblock{*cg} block and the aerodynamic tables. TAOS treats the value as
nondimensional.

A table can also be used:

\begin{lstlisting}[style=taosmanual]
*cg cg=(wt-cg) config=2
\end{lstlisting}

This requests table \tableid{wt-cg} to compute center of gravity. Like other tables,
it can depend on weight or on any permitted state or user-defined variable. The
example also assigns the user-defined variable \texttt{config}. A
\problemblock{*cg} block can assign values to any user-defined variables, but it is
normally used only for variables referenced by the center-of-gravity table.
